\documentclass[tikz,border=6pt]{standalone}
\usepackage[no-math]{fontspec}
\usepackage[brazil]{babel}
\usepackage{amsmath,amssymb}
\IfFileExists{C:/Windows/Fonts/Lato_Regular.ttf}{%
  \setmainfont{Lato}[
    Path = C:/Windows/Fonts/,
    UprightFont = Lato_Regular.ttf,
    BoldFont = Lato_Bold.ttf,
    ItalicFont = Lato_Italic.ttf,
    BoldItalicFont = Lato_BoldItalic.ttf
  ]
  \setsansfont{Lato}[
    Path = C:/Windows/Fonts/,
    UprightFont = Lato_Regular.ttf,
    BoldFont = Lato_Bold.ttf,
    ItalicFont = Lato_Italic.ttf,
    BoldItalicFont = Lato_BoldItalic.ttf
  ]
}{%
  \RequirePackage[default]{lato}
}
\renewcommand{\familydefault}{\sfdefault}

\usepackage{xcolor}
\usepackage{tikz}
\usepackage{pgfplots}
\pgfplotsset{compat=1.18}
\usetikzlibrary{arrows.meta,positioning,shapes.geometric,fit,backgrounds,calc,decorations.pathreplacing,patterns}

% Paleta institucional UEFS / CHF662
\definecolor{uefsNavy}{HTML}{1B2A4A}
\definecolor{uefsTeal}{HTML}{0E7C7B}
\definecolor{uefsForest}{HTML}{2D6A4F}
\definecolor{uefsGold}{HTML}{C4983B}
\definecolor{uefsGoldDark}{HTML}{8F6B1E}
\definecolor{uefsAlert}{HTML}{B02A37}
\definecolor{uefsOrange}{HTML}{D95F02}
\definecolor{uefsInk}{HTML}{22252A}
\definecolor{uefsMuted}{HTML}{5C6470}
\definecolor{uefsWater}{HTML}{2E86AB}
\definecolor{uefsSand}{HTML}{D9B382}
\definecolor{bgNavy}{HTML}{F2F5F9}
\definecolor{bgTeal}{HTML}{E9F5F4}
\definecolor{bgForest}{HTML}{EEF5F0}
\definecolor{bgGold}{HTML}{FCF7EB}
\definecolor{bgAlert}{HTML}{FDF2F2}
\definecolor{bgGray}{HTML}{F8FAFC}
\definecolor{lineGray}{HTML}{CFD8DC}

\tikzset{
  >=Stealth,
  aaCard/.style={draw=#1, fill=white, rounded corners=4pt, line width=1pt, inner sep=6pt, align=center, font=\small},
  aaBox/.style={draw=#1, rounded corners=4pt, line width=1.1pt, inner sep=6pt, align=center, font=\footnotesize\sffamily, text=uefsInk},
  aaTitle/.style={font=\normalsize\bfseries\sffamily, text=uefsNavy, align=center},
  aaSub/.style={font=\footnotesize\sffamily, text=uefsMuted, align=center},
  aaArrow/.style={->, >=Stealth, line width=1.4pt},
  aaDArrow/.style={<->, >=Stealth, line width=1.4pt}
}


\definecolor{ross1}{HTML}{D8F3DC}
\definecolor{ross2}{HTML}{95D5B2}
\definecolor{ross3}{HTML}{FFE066}
\definecolor{ross4}{HTML}{FFA94D}
\definecolor{ross5}{HTML}{E03131}

\begin{document}
\begin{tikzpicture}

% =========================================================================
% Painel (a): Ecodinâmica de Tricart (1977)
% =========================================================================
\begin{scope}
  \begin{axis}[
    width=7.1cm,
    height=6.7cm,
    ybar=4pt,
    bar width=16pt,
    title={(a) Ecodinâmica de Tricart (1977)\\Morfogênese vs. Pedogênese},
    title style={align=center, font=\small\bfseries\sffamily, text=uefsNavy, yshift=2pt},
    axis lines=left,
    axis line style={lineGray, line width=1pt},
    tick style={uefsMuted},
    ymin=0, ymax=112,
    ytick={0,20,40,60,80,100},
    yticklabel style={font=\scriptsize\sffamily, text=uefsInk},
    ylabel={\textbf{Predomínio relativo (\%)}},
    label style={font=\footnotesize\sffamily, text=uefsInk},
    symbolic x coords={Estáveis, Intergrades, Fort. Instáveis},
    xtick=data,
    xticklabels={{Meios\\Estáveis}, {Meios\\Intergrades}, {Meios Fort.\\Instáveis}},
    xticklabel style={align=center, font=\scriptsize\bfseries\sffamily, text=uefsInk},
    enlarge x limits=0.25,
    nodes near coords={\pgfmathprintnumber{\pgfplotspointmeta}\%},
    every node near coord/.append style={font=\tiny\bfseries\sffamily, text=uefsInk, yshift=1pt},
    legend style={
      at={(0.5,0.99)}, anchor=north,
      font=\scriptsize\sffamily, fill=white, fill opacity=0.95,
      draw=lineGray, rounded corners=3pt, inner sep=3pt,
      cells={anchor=west}
    }
  ]
    \addplot[fill=uefsForest, draw=none] coordinates {
      (Estáveis, 80) (Intergrades, 50) (Fort. Instáveis, 18)
    };
    \addlegendentry{Pedogênese (formação de solo)}

    \addplot[fill=uefsAlert, draw=none] coordinates {
      (Estáveis, 20) (Intergrades, 50) (Fort. Instáveis, 82)
    };
    \addlegendentry{Morfogênese (erosão / remoção)}
  \end{axis}
\end{scope}

% =========================================================================
% Painel (b): Matriz de Fragilidade Emergente (Ross, 1994)
% =========================================================================
\begin{scope}[xshift=8.65cm, yshift=0.15cm]
  \def\cw{1.42}
  \def\ch{1.00}

  \node[font=\small\bfseries\sffamily, text=uefsNavy] at ({2.5*\cw}, {5*\ch + 0.45})
    {(b) Matriz de Fragilidade Emergente (Ross, 1994)};

  % Desenho das 25 células (linha 1 embaixo até linha 5 em cima)
  \foreach \r/\c/\val/\bg/\fg/\lbl in {
    % Linha 1 (1 M. Fraca)
    0/0/1/ross1/uefsInk/M. Fraca,
    0/1/1/ross1/uefsInk/M. Fraca,
    0/2/2/ross2/uefsInk/Fraca,
    0/3/2/ross2/uefsInk/Fraca,
    0/4/3/ross3/uefsInk/Média,
    % Linha 2 (2 Fraca)
    1/0/1/ross1/uefsInk/M. Fraca,
    1/1/2/ross2/uefsInk/Fraca,
    1/2/2/ross2/uefsInk/Fraca,
    1/3/3/ross3/uefsInk/Média,
    1/4/4/ross4/uefsInk/Forte,
    % Linha 3 (3 Média)
    2/0/2/ross2/uefsInk/Fraca,
    2/1/2/ross2/uefsInk/Fraca,
    2/2/3/ross3/uefsInk/Média,
    2/3/4/ross4/uefsInk/Forte,
    2/4/4/ross4/uefsInk/Forte,
    % Linha 4 (4 Forte)
    3/0/2/ross2/uefsInk/Fraca,
    3/1/3/ross3/uefsInk/Média,
    3/2/4/ross4/uefsInk/Forte,
    3/3/4/ross4/uefsInk/Forte,
    3/4/5/ross5/white/M. Forte,
    % Linha 5 (5 M. Forte)
    4/0/3/ross3/uefsInk/Média,
    4/1/4/ross4/uefsInk/Forte,
    4/2/4/ross4/uefsInk/Forte,
    4/3/5/ross5/white/M. Forte,
    4/4/5/ross5/white/M. Forte%
  }{
    \filldraw[fill=\bg, draw=white, line width=1.4pt]
      ({\c*\cw}, {\r*\ch}) rectangle ({(\c+1)*\cw}, {(\r+1)*\ch});
    \node[align=center, text=\fg, font=\scriptsize\bfseries\sffamily]
      at ({(\c+0.5)*\cw}, {(\r+0.5)*\ch}) {\val\\[-1pt]\lbl};
  }

  % Moldura externa sutil
  \draw[lineGray, line width=1pt] (0,0) rectangle ({5*\cw}, {5*\ch});

  % Rótulos do eixo X (Uso e Cobertura)
  \foreach \c/\num/\txt in {
    0/1/Mata densa,
    1/2/Capoeira,
    2/3/Pastagem,
    3/4/Cultivo,
    4/5/Solo exp.%
  }{
    \draw[uefsMuted] ({(\c+0.5)*\cw}, 0) -- ({(\c+0.5)*\cw}, -0.09);
    \node[anchor=north, align=center, font=\scriptsize\sffamily, text=uefsInk]
      at ({(\c+0.5)*\cw}, -0.10) {\num\\[-1pt]\txt};
  }
  \node[font=\footnotesize\bfseries\sffamily, text=uefsInk]
    at ({2.5*\cw}, -0.95) {Classe de Proteção do Uso e Cobertura da Terra};

  % Rótulos do eixo Y (Fragilidade Potencial)
  \foreach \r/\txt in {
    0/1 M. Fraca,
    1/2 Fraca,
    2/3 Média,
    3/4 Forte,
    4/5 M. Forte%
  }{
    \draw[uefsMuted] (0, {(\r+0.5)*\ch}) -- (-0.09, {(\r+0.5)*\ch});
    \node[anchor=east, font=\scriptsize\sffamily, text=uefsInk]
      at (-0.11, {(\r+0.5)*\ch}) {\txt};
  }
  \node[rotate=90, font=\footnotesize\bfseries\sffamily, text=uefsInk]
    at (-1.65, {2.5*\ch}) {Fragilidade Potencial (Relevo + Solo + Clima)};
\end{scope}

\end{tikzpicture}
\end{document}
